\section{Principios fundamentales e ideas b\'asicas}
El handbook establece nueve principios fundamentales que aplican a las tareas y pr\'acticas de Requirements Engineering. Estos principios son la base metodol\'ogica para elicitar, documentar, validar y gestionar requisitos \cite[p.~16]{cpre}.

\begin{enumerate}[leftmargin=1.2em]
  \item \textbf{Orientaci\'on al valor.} Los requisitos son un medio para obtener resultados, no un fin documental en s\'i mismo.
  \item \textbf{Stakeholders.} RE busca satisfacer deseos y necesidades de las partes interesadas.
  \item \textbf{Entendimiento compartido.} Sin base com\'un entre actores, el desarrollo falla.
  \item \textbf{Contexto.} El sistema no se entiende aislado de su entorno.
  \item \textbf{Problema-requisito-soluci\'on.} Es una tr\'iada fuertemente entrelazada.
  \item \textbf{Validaci\'on.} Un requisito no validado no aporta valor real.
  \item \textbf{Evoluci\'on.} El cambio de requisitos es normal, no excepcional.
  \item \textbf{Innovaci\'on.} Repetir lo mismo no basta en muchos contextos.
  \item \textbf{Trabajo sistem\'atico y disciplinado.} Es imprescindible para mantener calidad y control.
\end{enumerate}

Como idea b\'asica transversal, el documento insiste en balancear coste y beneficio: la inversi\'on en RE reduce riesgo de retrabajo costoso en fases tard\'ias \cite[pp.~16--18]{cpre}.
